\subsection{トピックと参考文献の対応}
この表は、SWEBOKが示すトピック対参考資料の行列を、学習用に簡略化したものである。詳しく学ぶときは、各トピックに対応する文献を読む。
\par\smallskip\noindent\refstepcounter{table}\label{tab:ch18-15}表 \thetable: トピックと参考文献の対応の整理\par\smallskip
表 \thetable は、トピックと参考文献の対応の整理を比較・確認しやすい形で整理したものである。\par\smallskip
\begin{longtable}{>{\raggedright\arraybackslash}p{0.33\linewidth}>{\raggedright\arraybackslash}p{0.59\linewidth}}
\toprule
トピック & 主な参考資料 \\
\midrule
\endfirsthead
\toprule
トピック & 主な参考資料 \\
\midrule
\endhead
1 工学プロセス & Tockey, Return on Software \\
2 工学設計 & Voland, Engineering by Design; McConnell, Code Complete \\
3 抽象化とカプセル化 & McConnell, Code Complete \\
4 経験的方法と実験技法 & Montgomery and Runger, Applied Statistics and Probability for Engineers \\
5 統計分析 & Montgomery and Runger; Null and Lobur \\
6 モデリング・シミュレーション・プロトタイピング & Voland; Cheney and Kincaid; Sommerville; Fairley \\
7 測定 & Tockey; Voland; Fairley; Kan \\
8 標準 & Voland; Moore \\
9 根本原因分析 & Voland; Kan; Tockey; Ishikawa; Gano; Goldratt \\
10 インダストリー4.0とソフトウェア工学 & Nakagawa et al., Continuous Systems and Software Engineering for Industry 4.0 \\
\bottomrule
\end{longtable}
